
% ------------------------------------------------------------------------
% ------------------------------------------------------------------------
% abnTeX2: Modelo de Trabalho Academico (tese de doutorado, dissertacao de
% mestrado e trabalhos monograficos em geral) em conformidade com 
% ABNT NBR 14724:2011: Informacao e documentacao - Trabalhos academicos -
% Apresentacao
% ------------------------------------------------------------------------
% ------------------------------------------------------------------------

\documentclass[
	% -- opções da classe memoir --
	12pt,				% tamanho da fonte
	openright,			% capítulos começam em pág ímpar (insere página vazia caso preciso)
	oneside,				% para impressão em verso e anverso. Oposto a oneside
	a4paper,				% tamanho do papel. 
	% -- opções da classe abntex2 --
	chapter=TITLE,		% títulos de capítulos convertidos em letras maiúsculas
	%section=TITLE,			% títulos de seções convertidos em letras maiúsculas
	%subsection=TITLE,		% títulos de subseções convertidos em letras maiúsculas
	%subsubsection=TITLE,	% títulos de subsubseções convertidos em letras maiúsculas
	% -- opções do pacote babel --
	english,			% idioma adicional para hifenização
	french,			% idioma adicional para hifenização
	spanish,			% idioma adicional para hifenização
	brazil			% o último idioma é o principal do documento
	]{abntex2}

% ---
% Pacotes básicos 
% ---
\usepackage{lmodern}			% Usa a fonte Latin Modern			
\usepackage[T1]{fontenc}			% Selecao de codigos de fonte.
\usepackage[utf8]{inputenc}		% Codificacao do documento (conversão automática dos acentos)
\usepackage{lastpage}			% Usado pela Ficha catalográfica
\usepackage{indentfirst}			% Indenta o primeiro parágrafo de cada seção.
\usepackage[table,dvipsnames]{xcolor}				% Controle das cores
\usepackage{graphicx}				% para lidar com gráficos
\graphicspath{{./Figuras/}}
\DeclareGraphicsExtensions{.pdf,.jpeg,.jpg,.png}
\usepackage{microtype} 			% para melhorias de justificação
\usepackage{amsmath}
\usepackage{amsfonts}
\usepackage{bm}
\usepackage{enumitem}
\usepackage{siunitx}
\usepackage{multirow}
\usepackage{caption}
\usepackage[labelformat=simple]{subcaption}
\usepackage{xcolor}
\usepackage{booktabs}
\usepackage{pdfpages}
\usepackage{alphalph}
\usepackage{xfrac}
\usepackage{epigraph}
% ---

% ---
% Pacotes adicionais, usados apenas no âmbito do Modelo Canônico do abnteX2
% ---
\usepackage{lipsum}				% para geração de dummy text
% ---

% ---
% Pacotes de citações
% ---
\usepackage[brazilian,hyperpageref]{backref}	% Paginas com as citações na bibl
\usepackage[alf,abnt-etal-list=0,abnt-etal-cite=3,bibjustif,abnt-thesis-year=both]{abntex2cite}					% Citações padrão ABNT

\usepackage{IEEEtrantools}
%\begin{IEEEeqnarray}{rCl}
%	z(t) & = & 0  \\
%	t_i(t) & = & 0
%\end{IEEEeqnarray}

% *** ALIGNMENT PACKAGES ***
\usepackage{array}
\newcolumntype{L}[1]{>{\raggedright\let\newline\\\arraybackslash\hspace{0pt}}p{#1}}
\newcolumntype{C}[1]{>{\centering\let\newline\\\arraybackslash\hspace{0pt}}p{#1}}
\newcolumntype{R}[1]{>{\raggedleft\let\newline\\\arraybackslash\hspace{0pt}}p{#1}}

% --- 
% CONFIGURAÇÕES DE PACOTES
% --- 

% ---
% Configurações do pacote backref
% Usado sem a opção hyperpageref de backref
\renewcommand{\backrefpagesname}{Citado na(s) página(s):~}
% Texto padrão antes do número das páginas
\renewcommand{\backref}{}
% Define os textos da citação
\renewcommand*{\backrefalt}[4]{
	\ifcase #1 %
		Nenhuma citação no texto.%
	\or
		Citado na página #2.%
	\else
		Citado #1 vezes nas páginas #2.%
	\fi}%
% ---

% ---
% Configurações do pacote siunitx
\sisetup{
   output-decimal-marker = {,} ,
   group-separator = {.}
}
% ---
% Configurações do pacote subcaption
%\renewcommand\thesubfigure{(\alph{subfigure})}
\renewcommand*{\thesubfigure}{(\alphalph{\value{subfigure}})}
% ---

% Configurações do pacote epigraph
\setlength{\epigraphrule}{0pt}

% ---
% Definições de comandos
% ---
\newcommand{\myVec}[1]{\bm{#1}}
\newcommand{\myHat}[1]{\hat{\bm{#1}}}
\newcommand{\myMatrix}[1]{\bm{\mathit{#1}}}
\newcommand{\ud}{\,\mathrm{d}}
\newcommand{\tabLabel}[1]{\multicolumn{1}{c}{\textbf{#1}}}
\newcommand{\tabLabelMrow}[2]{\multicolumn{1}{c}{\multirow{#1}{*}{\textbf{#2}}}}
\newcommand{\rowSpace}[1]{\renewcommand{\arraystretch}{#1}}
\newcommand{\minitab}[2][c]{\begin{tabular}{#1}#2\end{tabular}}

\newlist{sumariotese}{enumerate}{4}
\setlist[sumariotese,1]{label=\arabic*}
\setlist[sumariotese,2]{label=\arabic{sumariotesei}.\arabic*}
\setlist[sumariotese,3]{label=\arabic{sumariotesei}.\arabic{sumarioteseii}.\arabic*}
\setlist[sumariotese,4]{label=\arabic{sumariotesei}.\arabic{sumarioteseii}.\arabic{sumarioteseiii}.\arabic*}

% ---
% Informações de dados para CAPA e FOLHA DE ROSTO
% ---
\titulo{Título do projeto}
\autor{Nome do Aluno}
\local{Manaus}
\data{2021}
\orientador{Prof. Dr. Danilo de Santana Chui}
%\coorientador{Equipe \abnTeX}
\instituicao{Faculdade de Tecnologia, Universidade Federal do Amazonas}
\tipotrabalho{Relatório de PIBIC}
% O preambulo deve conter o tipo do trabalho, o objetivo, 
% o nome da instituição e a área de concentração 
\preambulo{Relatório de Projeto de Iniciação Científica para cumprir com os requisitos estabelecidos no Edital 081/2019-PROPESP/UFAM do Programa Institucional de Bolsas de Iniciação Científica da Universidade Federal do Amazonas.}
% ---


% ---
% Configurações de aparência do PDF final

% alterando o aspecto da cor azul
\definecolor{blue}{RGB}{41,5,195}

% informações do PDF
\makeatletter
\hypersetup{
     	%pagebackref=true,
	pdftitle={\@title}, 
	pdfauthor={\@author},
    	pdfsubject={\imprimirpreambulo},
	pdfcreator={LaTeX with abnTeX2},
	pdfkeywords={ufam}{ft}{engenharia mecânica}{mecânica do contínuo}, 
	colorlinks=true,       		% false: boxed links; true: colored links
    	linkcolor=black,          	% color of internal links
    	citecolor=black,        		% color of links to bibliography
    	filecolor=black,      		% color of file links
	urlcolor=black,
	bookmarksdepth=4
}
\makeatother
% --- 

% --- 
% Espaçamentos entre linhas e parágrafos 
% --- 

% O tamanho do parágrafo é dado por:
\setlength{\parindent}{1.3cm}

% Controle do espaçamento entre um parágrafo e outro:
\setlength{\parskip}{0.2cm}  % tente também \onelineskip

% Espessura das bordas das tabelas
\setlength\heavyrulewidth{0.12em}

% ---
% compila o indice
% ---
\makeindex
% ---

% NEW COMMANDS
\providecommand{\en}[1]{\ensuremath{\textrm{\textsc{e}-}{#1}}}
\providecommand{\ep}[1]{\ensuremath{\textrm{\textsc{e}}{#1}}}

\renewcommand{\imprimircapa}{
	\begin{capa}
		\begin{center}
		\begin{tabular}{L{0.13\textwidth}C{0.68\textwidth}R{0.14\textwidth}}
			\multirow{3}{*}{\includegraphics[height=2cm]{ufam.png}} & 
			{\ABNTEXchapterfont\large UNIVERSIDADE FEDERAL DO AMAZONAS} \vspace*{0.2cm} &
			\multirow{3}{*}{\includegraphics[height=2.2cm]{engmec.png}} \\
			& {\ABNTEXchapterfont\large FACULDADE DE TECNOLOGIA} \vspace*{0.2cm} & \\
			& {\ABNTEXchapterfont\large DEPTO. DE ENGENHARIA MECÂNICA} & \\
		\end{tabular}
		\\[0.5cm]
		\includegraphics[width=3cm]{propesp_logo.png} \\
		{\ABNTEXchapterfont\large Pró-Reitoria de Pesquisa e Pós Graduação} \\
		{\ABNTEXchapterfont\large Programa Institucional de Bolsas de Iniciação Científica (PIBIC)} \\
		{\ABNTEXchapterfont\large Programa de Apoio à Iniciação Científica (PAIC)} \\[2.5cm]
		{\ABNTEXchapterfont\bfseries\LARGE\MakeUppercase RELATÓRIO PARCIAL} \\[0.5cm]
		{\ABNTEXchapterfont\bfseries\LARGE\MakeUppercase\imprimirtitulo} \\[4.5cm]
		{\ABNTEXchapterfont\bfseries\LARGE\MakeUppercase\imprimirautor}
		\vfill
		{\ABNTEXchapterfont\large\imprimirlocal} \\[0.2cm]
		{\ABNTEXchapterfont\large\imprimirdata}
		\vspace*{1cm}
		\end{center}
	\end{capa}
}

\makeatletter
\renewcommand{\folhaderostocontent}{
	\begin{center}
		{\ABNTEXchapterfont\large\imprimirautor} \\[6cm]
		\begin{center}
			{\ABNTEXchapterfont\bfseries\Large RELATÓRIO PARCIAL} \\[1cm]
			{\ABNTEXchapterfont\bfseries\Large\imprimirtitulo} \\
		\end{center}
		\vfill
		\abntex@ifnotempty{\imprimirpreambulo}{
			\hspace{4cm}
			\begin{minipage}{0.68\textwidth}
				\SingleSpacing
				\imprimirpreambulo \\[0.5cm]
			\end{minipage}
		}
		\vfill
		{\large\imprimirlocal}
		\par
		{\large\imprimirdata}
		\vspace*{1cm}
	\end{center}
}
\makeatother

% TEOREMAS
\newtheorem{teorema}{Teorema}[chapter]
\newtheorem{definicao}{Definição}[chapter]
\newtheorem{algoritmo}{Algoritmo}[chapter]

% OPERADORES MATEMATICOS
\DeclareMathOperator{\cov}{cov}

% ----
% Início do documento
% ----
\begin{document}

% Retira espaço extra obsoleto entre as frases.
\frenchspacing 

% ----------------------------------------------------------
% ELEMENTOS PRÉ-TEXTUAIS
% ----------------------------------------------------------
% \pretextual

% ---
% Capa
% ---
\imprimircapa
% ---

% ---
% Folha de rosto
% (o * indica que haverá a ficha bibliográfica)
% ---
%\imprimirfolhaderosto*
\imprimirfolhaderosto*
% ---

% ---
% Inserir a ficha bibliografica
% ---

% Isto é um exemplo de Ficha Catalográfica, ou ``Dados internacionais de
% catalogação-na-publicação''. Você pode utilizar este modelo como referência. 
% Porém, provavelmente a biblioteca da sua universidade lhe fornecerá um PDF
% com a ficha catalográfica definitiva após a defesa do trabalho. Quando estiver
% com o documento, salve-o como PDF no diretório do seu projeto e substitua todo
% o conteúdo de implementação deste arquivo pelo comando abaixo:
%
% \begin{fichacatalografica}
%     \includepdf{EPUSP-Catalogacao-na-Fonte_corrigida_passo2.pdf}
% \end{fichacatalografica}
%
\begin{fichacatalografica}
	\begin{minipage}[c]{12.5cm}	
	\noindent Autorizo a reprodução e divulgação total ou parcial deste trabalho, por qualquer meio convencional ou eletrônico, para fins de referência, estudo ou pesquisa, desde que citada a fonte. \\[1cm]
		\begin{center}
			Catalogação na publicação \\
			Serviço de Biblioteca e Documentação \\
			Faculdade de Tecnologia da Universidade Federal do Amazonas
		\end{center}
	\end{minipage}
	
	\vspace*{\fill}					% Posição vertical
	\hrule							% Linha horizontal
	\begin{center}					% Minipage Centralizado
	\begin{minipage}[c]{12.5cm}		% Largura
	
	\imprimirautor
	
	\hspace{0.5cm} \imprimirtitulo  / \imprimirautor. --
	\imprimirlocal, \imprimirdata-
	
	\hspace{0.5cm} \pageref{LastPage} p. : il. (algumas color.) ; 30 cm.\\
	
	\hspace{0.5cm} \imprimirorientadorRotulo~\imprimirorientador\\
	
	\hspace{0.5cm}
	\parbox[t]{\textwidth}{\imprimirtipotrabalho~--~\imprimirinstituicao,
	\imprimirdata.}\\
	
	\hspace{0.5cm}
		1. PIBIC.
		2. Engenharia Mecânica.
		I. Universidade Federal do Amazonas.
		II. Faculdade de Tecnologia.\\ 			
	
	\hspace{8.75cm} CDU 02:141:005.7\\
	
	\end{minipage}
	\end{center}
	\hrule
\end{fichacatalografica}
\cleardoublepage
% ---

% ---
% Inserir errata
% ---
%\begin{errata}
%\noindent CHUI, D. S. \textbf{Controle de chama através do monitoramento da qualidade de nebulização por visão computacional e filtro de Kalman}. 2016. XXX f. Tese (Doutorado) - Departamento de Engenharia Mecânica, Escola Politécnica, Universidade de São Paulo, São Paulo, 2016.
%
%\begin{table}[htb]
%\center
%\footnotesize
%\begin{tabular}{|p{1.5cm}|p{1.5cm}|p{5cm}|p{5cm}|}
%  \hline
%   \textbf{Folha} & \textbf{Linha}  & \textbf{Onde se lê}  & \textbf{Leia-se}  \\
%    \hline
%    1 & 10 & auto-conclavo & autoconclavo\\
%   \hline
%\end{tabular}
%\end{table}
%
%\end{errata}
% ---

% ---
% Inserir folha de aprovação
% ---

% Isto é um exemplo de Folha de aprovação, elemento obrigatório da NBR
% 14724/2011 (seção 4.2.1.3). Você pode utilizar este modelo até a aprovação
% do trabalho. Após isso, substitua todo o conteúdo deste arquivo por uma
% imagem da página assinada pela banca com o comando abaixo:
%
% \includepdf{folhadeaprovacao_final.pdf}
%
%\begin{folhadeaprovacao}
 
%    \noindent Nome: \imprimirautor \\[0.5cm]
%	Título: \imprimirtitulo
%
%	\hspace{4cm}
%	\begin{minipage}{0.66\textwidth}
%		\SingleSpacing
%		\imprimirpreambulo \\[0.5cm]
%		Área de Concentração: Engenharia de Controle e Automação Mecânica. \\[0.2cm]
%	\end{minipage}
%	\noindent Aprovado em: 17/02/2020\\
	
%\end{folhadeaprovacao}
% ---

% ---
% Dedicatória
% ---
%\include{dedicatoria}
% ---

% ---
% Agradecimentos
% ---
%\include{agradecimentos}
% ---

% ---
% Epígrafe
% ---
%\include{epigrafe}
% ---

\chapter*[Dados do Projeto]{Dados do Projeto}
\begin{center}
    \noindent {\ABNTEXchapterfont
    \begin{tabular}{|llll|}
        \hline
        \rowcolor{CornflowerBlue} \multicolumn{4}{|c|}{\textbf{EDIÇÃO:} PIBIC/PAIC 2020/2021} \\
        \hline
        \rowcolor{CornflowerBlue} \multicolumn{4}{|l|}{\textbf{RECURSOS HUMANOS}} \\
        \hline
        \multicolumn{2}{|l|}{\textbf{Nome do orientador:}} & \multicolumn{2}{|l|}{\textbf{Bolsa:}} \\
        \multicolumn{2}{|l|}{\imprimirorientador} & \multicolumn{2}{|l|}{\multirow{3}{*}{
            \begin{tabular}{ll}
             ( ) CNPQ & ( ) UFAM   \\
             ( ) FAPEAM & ( ) VOLUNTÁRIO  
            \end{tabular}
        }} \\
        \cline{1-2}
        \multicolumn{2}{|l|}{\textbf{Nome do aluno:}} & \multicolumn{2}{|l|}{} \\
        \multicolumn{2}{|l|}{\imprimirautor} & \multicolumn{2}{|l|}{} \\
        \hline
        \rowcolor{CornflowerBlue} \multicolumn{4}{|l|}{\textbf{IDENTIFICAÇÃO DO PROJETO}} \\
        \hline
        \multicolumn{3}{|l|}{\textbf{Título:}} & \multicolumn{1}{|l|}{\textbf{Código do Projeto:}}\\
        \multicolumn{3}{|l|}{\imprimirtitulo} & \multicolumn{1}{|l|}{PIB-ENG-XXXX}\\
        \hline
        \multicolumn{4}{|l|}{\textbf{Área de Conhecimento:}} \\
        ( ) Agrárias & \multicolumn{2}{l}{\textbf{(X) Engenharias}} & ( ) Multidisciplinar \\
        ( ) Biológicas & \multicolumn{2}{l}{( ) Exatas e da Terra} & ( ) Saúde \\
        ( ) Ciências Humanas & \multicolumn{2}{l}{( ) Linguística, Letras e Artes} & ( ) Sociais Aplicadas \\
        \hline
        \rowcolor{CornflowerBlue} \multicolumn{4}{|l|}{\textbf{COMITÊ DE ÉTICA EM PESQUISA COM HUMANOS (CEP)}} \\
        \rowcolor{CornflowerBlue} \multicolumn{4}{|l|}{\textbf{OU ANIMAIS (CEUA)}} \\
        \hline
        \multicolumn{4}{|l|}{( ) Aprovado - número do protocolo:} \\
        \multicolumn{4}{|l|}{\textbf{(X) Não se aplica}} \\
        \multicolumn{4}{|l|}{Caso o projeto não esteja aprovado, justifique:} \\
        \multicolumn{4}{|l|}{} \\
        \multicolumn{4}{|l|}{} \\
        \multicolumn{4}{|l|}{} \\
        \hline
    \end{tabular}}  
\end{center}
\cleardoublepage

% ---
% RESUMOS
% ---

% resumo em português
\setlength{\absparsep}{18pt} % ajusta o espaçamento dos parágrafos do resumo
\begin{resumo}
	\noindent NOME DO AUTOR {\textit\imprimirtitulo.} \imprimirdata. \pageref{LastPage} p. \imprimirtipotrabalho~--~\imprimirinstituicao, \imprimirlocal, \imprimirdata. \\[0.5cm]
	
	(Aqui vai o resumo)
	
	\textbf{Palavras-chaves}: 
\end{resumo}

% resumo em inglês
%\begin{resumo}[Abstract]
%	\begin{otherlanguage*}{english}
%
%   	\textbf{Keywords}: 
%	\end{otherlanguage*}
%\end{resumo}
%\cleardoublepage

%% resumo em francês 
%\begin{resumo}[Résumé]
% \begin{otherlanguage*}{french}
%    Il s'agit d'un résumé en français.
% 
%   \textbf{Mots-clés}: latex. abntex. publication de textes.
% \end{otherlanguage*}
%\end{resumo}
%
%% resumo em espanhol
%\begin{resumo}[Resumen]
% \begin{otherlanguage*}{spanish}
%   Este es el resumen en español.
%  
%   \textbf{Palabras clave}: latex. abntex. publicación de textos.
% \end{otherlanguage*}
%\end{resumo}
%% ---

% ---
% inserir lista de ilustrações
% ---
\pdfbookmark[0]{\listfigurename}{lof}
\listoffigures*
\cleardoublepage
% ---

% ---
% inserir lista de tabelas
% ---
\pdfbookmark[0]{\listtablename}{lot}
\listoftables*
\cleardoublepage
% ---

% ---
% inserir lista de abreviaturas e siglas
% ---
\begin{siglas}
	\item[UFAM] Universidade Federal do Amazonas.
\end{siglas}
% ---

% ---
% inserir lista de símbolos
% ---
%\begin{simbolos}
%	\item[$A$] área da imagem $I^{pb}$.
%\end{simbolos}
% ---

% ---
% inserir o sumario
% ---
\pdfbookmark[0]{\contentsname}{toc}
\tableofcontents*
\cleardoublepage
% ---



% ----------------------------------------------------------
% ELEMENTOS TEXTUAIS
% ----------------------------------------------------------
\textual

\include{1_introducao}

\include{2_revisao_literatura}

\include{3_metodologia}

\include{4_resultados}

\include{5_conclusoes}

%% ----------------------------------------------------------
%% Finaliza a parte no bookmark do PDF
%% para que se inicie o bookmark na raiz
%% e adiciona espaço de parte no Sumário
%% ----------------------------------------------------------
%\phantompart
%
%% ---
%% Conclusão (outro exemplo de capítulo sem numeração e presente no sumário)
%% ---
%\chapter[Conclusão]{Conclusão}
%%\addcontentsline{toc}{chapter}{Conclusão}
%% ---
%
%\lipsum[31-33]

% ----------------------------------------------------------
% ELEMENTOS PÓS-TEXTUAIS
% ----------------------------------------------------------
\postextual
% ----------------------------------------------------------

% ----------------------------------------------------------
% Referências bibliográficas
% ----------------------------------------------------------
\bibliography{referencias}
% ----------------------------------------------------------
% Glossário
% ----------------------------------------------------------
%
% Consulte o manual da classe abntex2 para orientações sobre o glossário.
%
%\glossary

% ----------------------------------------------------------
% Apêndices
% ----------------------------------------------------------

% ---
% Inicia os apêndices
% ---
\begin{apendicesenv}

% Imprime uma página indicando o início dos apêndices
% \partapendices

%\include{apendice_imagens_proc}

\end{apendicesenv}
% ---


% ----------------------------------------------------------
% Anexos
% ----------------------------------------------------------

%% ---
%% Inicia os anexos
%% ---
%\begin{anexosenv}
%
%% Imprime uma página indicando o início dos anexos
%\partanexos
%
%% ---
%\chapter{Morbi ultrices rutrum lorem.} 
%% ---
%\lipsum[30]
%
%\end{anexosenv}

%---------------------------------------------------------------------
% INDICE REMISSIVO
%---------------------------------------------------------------------
\phantompart
\printindex
%---------------------------------------------------------------------

\end{document}

